\chapter{Mortgages and Agencies}\label{ch:m2:mortgages-and-agencies}

An American who borrows to buy a house usually borrows for thirty years at a
fixed rate and keeps the right to repay at any time, without penalty. When
rates fall, millions of borrowers refinance, and the investors who own their
mortgages get their money back just when they can reinvest it only at the lower
rate. The mortgage bond that should have gained from the fall gains little; its
duration collapses; and the investors who hedged it must buy duration in a
hurry, from Treasuries or swaps. In March 2003 two Federal Reserve economists
estimated that a fall of half a point in rates would have shortened the duration
of outstanding mortgage securities by as much as receiving fixed on USD~289
billion of new ten-year swaps. This chapter explains the mortgage-backed
security that makes the American fixed-rate mortgage possible, the forward
market in which almost all of it trades, and the \omterm{def:m2:mortgages-and-agencies:convexity}{negative convexity} that makes
its holders move the rest of the rates market.

\section{The agency pass-through}

\begin{definition}[Agency mortgage-backed security, pass-through]\label{def:m2:mortgages-and-agencies:mbs}
An \emph{agency mortgage-backed security}\index{agency mortgage-backed security}
is a security backed by a pool of residential mortgages whose timely payment of
interest and principal is guaranteed by a government-sponsored enterprise
(Fannie Mae, Freddie Mac) or by a government agency (Ginnie Mae, whose guarantee
carries the full faith and credit of the United States). The basic form is the
\emph{pass-through}\index{pass-through}: each month it passes the borrowers'
payments, interest at the pool's coupon and all principal, scheduled and
prepaid, to the investors pro rata, after the servicer and the guarantor have
taken their fees.
\end{definition}

The pool's loans pay a gross rate, the weighted average coupon; investors
receive a lower net coupon, the difference paying the servicer who collects the
payments and the guarantor who bears the credit risk. A pool of 6.5\% loans may
back a 6\% \omterm{def:m2:mortgages-and-agencies:mbs}{pass-through}. The investor has, in effect, no credit risk and one
large risk instead: the timing of the principal.

\begin{omfigure}
\centering
\begin{tikzpicture}[font=\small,
  box/.style={draw,thick,rounded corners=2pt,align=center,minimum height=1cm,minimum width=2.2cm,inner sep=3pt},
  arr/.style={-{Stealth[length=2.5mm]},thick}]
\node[box,draw=omDef,fill=omDef!8] (b) at (0,0) {borrowers\\(6.5\% loans)};
\node[box,draw=omThm,fill=omThm!8] (s) at (4.3,0) {servicer};
\node[box,draw=omThm,fill=omThm!8] (g) at (8.6,0) {pool with\\agency guarantee};
\node[box,draw=omProp,fill=omProp!8] (i) at (12.9,0) {investors\\(6\% coupon)};
\draw[arr] (b) -- node[above=6pt,font=\footnotesize,align=center] {payments,\\prepayments} (s);
\draw[arr] (s) -- node[above=6pt,font=\footnotesize] {net of fee} (g);
\draw[arr] (g) -- node[above=6pt,font=\footnotesize] {pro rata} (i);
\node[font=\footnotesize,text=omThm,align=center] at (4.3,-1.2) {keeps a servicing fee};
\node[font=\footnotesize,text=omThm,align=center] at (8.6,-1.25) {keeps a guarantee fee,\\makes investors whole on default};
\end{tikzpicture}

\omcaption{How a \omterm{def:m2:mortgages-and-agencies:mbs}{pass-through} works. Borrowers pay the servicer; the servicer
and the guarantor keep their fees; investors receive interest at the net
coupon and every dollar of principal, including \omterm{def:m2:mortgages-and-agencies:cpr}{prepayments}, in proportion to
their holdings. Schematic, with the chapter's illustrative rates.}\label{fig:m2:mortgages-and-agencies:passthrough}
\end{omfigure}

The market is among the largest in fixed income: a 2022 survey by economists
of the Federal Reserve Bank of New York put American mortgage-backed securities
above USD~11 trillion outstanding and their trading near USD~300 billion a day,
most of it in the agency securities. In mid-2021 banks held about a third of
them and the Federal Reserve, the largest single holder, about a quarter,
bought in its programmes of asset purchases after 2008 and after 2020
(\cref{fig:m2:mortgages-and-agencies:fed}).

\begin{dated}{2026-09}{Agency mortgages}\label{dat:m2:mortgages-and-agencies:market}
Agency MBS trading averaged USD~367.0 billion a day in 2026 to August (SIFMA);
MBS issuance was USD~1\,432.5 billion over the same period. The Federal Reserve
held USD~1\,913.5 billion of MBS on 16 September 2026, down from a peak of
USD~2\,740.2 billion on 13 April 2022, letting its holdings run off. The
thirty-year fixed mortgage rate in the Freddie Mac survey was 6.95\% on 17
September 2026.
\end{dated}

\section{Prepayment}

\begin{definition}[Prepayment, conditional prepayment rate, PSA benchmark]\label{def:m2:mortgages-and-agencies:cpr}
A \emph{prepayment}\index{prepayment} is a repayment of mortgage principal
ahead of the amortisation schedule: a sale of the house, a refinancing, a
partial repayment, or a default bought out of the pool by the guarantor. The
\emph{conditional prepayment rate}\index{conditional prepayment rate} (CPR) is
the annualised fraction of the balance, after scheduled principal, prepaid in a
month; the monthly fraction, the single monthly mortality, is
$\mathrm{SMM} = 1 - (1-\mathrm{CPR})^{1/12}$. The \emph{PSA
benchmark}\index{PSA benchmark} is a standard speed path: 100\% PSA is a CPR of
0.2\% in a loan's first month, rising by 0.2\% a month to 6\% in its thirtieth
and constant after; 200\% PSA doubles every rate.
\end{definition}

\omterm{def:m2:mortgages-and-agencies:cpr}{Prepayment} speeds follow the borrowers' incentives. New loans prepay slowly,
since few people move or refinance in the first months; this is the ramp the
\omterm{def:m2:mortgages-and-agencies:cpr}{PSA benchmark} describes (\cref{fig:m2:mortgages-and-agencies:psa}). The
benchmark is a quoting convention, not a model: the Ginnie Mae offering
documents that define it say that it does not describe history or predict
anything. What drives speeds is the \emph{refinancing incentive}, the gap
between the rate borrowers pay and the rate they could get today. Out of the
money, a pool prepays at the slow pace of house sales; in the money, several
times faster; the transition between the two, an S-shaped curve in the
incentive, is where the risk lies. The chapter's illustrative curve runs from
a CPR of 6\% far out of the money to 50\% far in it, and gives 11.9\% at the
money.

\begin{omfigure}
\begin{tikzpicture}
\begin{axis}[width=11cm,height=4.8cm,
  xlabel={loan age (months)},ylabel={CPR (\%)},
  tick label style={font=\small},label style={font=\small},
  xmin=0,xmax=60,ymin=0,ymax=13,grid=major,grid style={gray!25},
  legend columns=3,legend style={font=\footnotesize,at={(0.5,-0.32)},anchor=north,draw=none,fill=none,/tikz/every even column/.append style={column sep=8pt}},
  table/col sep=comma]
\addplot[omDef,thick,dashed] table[x=age,y=psa50]{figdata/markets-2/12-mortgages-and-agencies/psa.csv};
\addplot[omThm,very thick] table[x=age,y=psa100]{figdata/markets-2/12-mortgages-and-agencies/psa.csv};
\addplot[omProp,thick,dotted] table[x=age,y=psa200]{figdata/markets-2/12-mortgages-and-agencies/psa.csv};
\legend{50\% PSA,100\% PSA,200\% PSA}
\end{axis}
\end{tikzpicture}

\omcaption{The \omterm{def:m2:mortgages-and-agencies:cpr}{PSA benchmark}: annual \omterm{def:m2:mortgages-and-agencies:cpr}{prepayment} rate by loan age, rising
linearly for thirty months and flat after. Speeds are quoted as multiples of
this path. Data: the chapter's tutorial, from the definition in Ginnie Mae's
offering documents.}\label{fig:m2:mortgages-and-agencies:psa}
\end{omfigure}

\begin{example}[A new pool]\label{ex:m2:mortgages-and-agencies:pool}
USD~100 million of new thirty-year 6.5\% loans back a 6\% \omterm{def:m2:mortgages-and-agencies:mbs}{pass-through}. The
level monthly payment is USD~632\,068; in the first month it contains
USD~541\,667 of interest and USD~90\,401 of scheduled principal, and at 100\%
PSA (a CPR of 0.2\%) USD~16\,667 is prepaid. Investors receive USD~500\,000 of
interest, the 6\% coupon, and all USD~107\,068 of principal. The weighted
average life, the average time to the return of a dollar of principal, is 19.6
years with no \omterm{def:m2:mortgages-and-agencies:cpr}{prepayment}, 11.5 at 100\% PSA and 5.8 at 300\%.
\end{example}

\section{The to-be-announced market and dollar rolls}

\begin{definition}[To-be-announced trade, specified pool, dollar roll]\label{def:m2:mortgages-and-agencies:tba}
A \emph{to-be-announced trade}\index{to-be-announced trade} (TBA) is a forward
trade in agency \omterm{def:m2:mortgages-and-agencies:mbs}{pass-throughs} in which the buyer and the seller agree on six
parameters (agency, coupon, maturity, price, face value and settlement month),
but not on the pools: two business days before the monthly settlement date the
seller names pools that meet the industry's good-delivery guidelines. A
\emph{specified pool}\index{specified pool} is a pool traded by its
identifier, at a price above the TBA price (a pay-up) when its loans are
expected to prepay more slowly. A \emph{dollar roll}\index{dollar roll} is the
sale of a TBA for one settlement month together with the purchase of the same
TBA for a later month.
\end{definition}

The TBA market turns around a million distinct pools into a handful of liquid
contracts, one per agency family, coupon and settlement month; since mid-2019
pools of Fannie Mae and Freddie Mac trade in the same Uniform MBS contracts. The
seller delivers the least valuable pools that qualify, so the TBA works like a
futures contract with a \omterm{def:m2:bond-futures:ctd}{cheapest-to-deliver} (\cref{ch:m2:bond-futures}); pools
whose borrowers are expected to prepay more slowly are worth more, are kept out
of TBA delivery and are sold as \omterm{def:m2:mortgages-and-agencies:tba}{specified pools}. The concentration pays: the 2022 survey reports about USD~261
billion of the agency market's USD~288 billion daily trading in TBAs, at an
estimated one-way cost of about one basis point, against about forty for
\omterm{def:m2:mortgages-and-agencies:tba}{specified pools}. Lenders use the market to hedge the loans they have promised
but not yet made, by selling TBAs forward.

A \omterm{def:m2:mortgages-and-agencies:tba}{dollar roll} is the market's financing trade. The roll seller, long TBAs for
the front month, sells them and buys them back for the next month: she gives
up a month of coupon and principal paydown, and keeps the cash for a month.
The difference between the two prices, the \emph{drop}, sets the rate at which
she has in effect borrowed.

\begin{proposition}[Implied financing rate of a dollar roll]\label{prop:m2:mortgages-and-agencies:roll}
Sell the front month at $P_1$ and buy the back month at $P_2$, $\tau$ days
later, on a \omterm{def:m2:mortgages-and-agencies:mbs}{pass-through} of coupon $c$ whose principal pays down a fraction
$s$ over the month (returned at par). The roll's implied financing rate $r$
solves
\[
P_1\left(1 + r\,\frac{\tau}{360}\right) \;=\; (1-s)\,P_2 + 100\,s + 100\,\frac{c}{12}.
\]
\end{proposition}

\begin{proof}
Holding the securities for the month ends with face $100(1-s)$, worth
$(1-s)P_2$, plus the paydown $100s$ and the coupon $100c/12$. Rolling
instead yields $P_1$ invested for $\tau$ days, from which the same face can be
bought back at $P_2$. The two positions end with the same securities when the
cash sides match, which is the equation. Accrued interest, equal at both
settlement dates for a day-of-month convention, cancels.
\end{proof}

\begin{example}[A roll]\label{ex:m2:mortgages-and-agencies:roll}
The 6\% TBA trades at 100.50 for the front month and 100.25 for the next, a
drop of a quarter point; the pool pays down 1.16\% a month. The roll finances at
2.95\% a year. With repo at 4\%, holding the securities and financing them in
repo costs more than rolling: the roll is \emph{special}, as the dealers say,
and the drop that would make the two equivalent is 0.16. A roll trades special
when securities for the front month are scarce, for instance when a large buyer
is taking delivery of many pools; the Federal Reserve has itself used rolls in
its operations, including in 2020.
\end{example}

\section{Negative convexity and its hedging flows}

\begin{definition}[Effective duration, negative convexity]\label{def:m2:mortgages-and-agencies:convexity}
The \emph{effective duration}\index{effective duration} of a security whose
cash flows depend on rates is $(P_- - P_+)/(2P\,\Delta y)$, where $P_\pm$ are
its prices after shifting rates by $\pm\Delta y$ and re-running the model of
its cash flows. A security has \emph{negative convexity}\index{negative
convexity} over a range of rates when its effective duration falls as rates
fall: its price rises less than it falls for equal moves.
\end{definition}

An ordinary bond has positive convexity (\cref{ch:m2:government-bonds}): as
rates fall its duration lengthens and its price accelerates. A \omterm{def:m2:mortgages-and-agencies:mbs}{pass-through} does
the opposite near the money, because every borrower holds a call option on the
loan. As rates fall the option moves into the money, \omterm{def:m2:mortgages-and-agencies:cpr}{prepayments} accelerate,
principal comes back at par and the price is capped near par plus a little.
\Cref{fig:m2:mortgages-and-agencies:priceyield} shows it on the chapter's pool:
the \omterm{def:m2:mortgages-and-agencies:mbs}{pass-through} gains 2.07 points on a 100-basis-point fall and loses 6.11 on
a 100-basis-point rise, while the same pool with its speed frozen would gain
5.04 and lose 4.60.

\begin{omfigure}
\begin{tikzpicture}
\begin{axis}[width=11cm,height=5cm,
  xlabel={discount yield (\%)},ylabel={price},
  tick label style={font=\small},label style={font=\small},
  xmin=3.9,xmax=8.1,ymin=86,ymax=112,grid=major,grid style={gray!25},
  legend columns=2,legend style={font=\footnotesize,at={(0.5,-0.3)},anchor=north,draw=none,fill=none,/tikz/every even column/.append style={column sep=8pt}},
  table/col sep=comma]
\addplot[omThm,very thick,mark=*,mark size=1.2pt] table[x=y,y=mbs]{figdata/markets-2/12-mortgages-and-agencies/priceyield.csv};
\addplot[omDef,thick,dashed] table[x=y,y=static]{figdata/markets-2/12-mortgages-and-agencies/priceyield.csv};
\draw[gray] (axis cs:6,86) -- (axis cs:6,112);
\node[font=\footnotesize,anchor=west] at (axis cs:6.08,109) {at the money};
\legend{pass-through (speeds respond to rates),same pool with speeds frozen}
\end{axis}
\end{tikzpicture}

\omcaption{Price of the chapter's 6\% \omterm{def:m2:mortgages-and-agencies:mbs}{pass-through} against the discount yield,
when \omterm{def:m2:mortgages-and-agencies:cpr}{prepayment} speeds follow the refinancing S-curve and when they are frozen
at their at-the-money value. Below the money the \omterm{def:m2:mortgages-and-agencies:mbs}{pass-through}'s price flattens
out: \omterm{def:m2:mortgages-and-agencies:convexity}{negative convexity}. Illustrative model; data: the chapter's
tutorial.}\label{fig:m2:mortgages-and-agencies:priceyield}
\end{omfigure}

\begin{proposition}[The convexity hedge]\label{prop:m2:mortgages-and-agencies:hedge}
A holder who keeps a \omterm{def:m2:mortgages-and-agencies:mbs}{pass-through} portfolio of face $F$ hedged to zero
duration must, after a rate move from $y_0$ to $y_1$, change its hedge by
\[
\Delta\mathrm{DV01} \;=\; \frac{F}{100}\,\bigl(P(y_1) D(y_1) - P(y_0) D(y_0)\bigr)\times 10^{-4},
\]
with $P$ the price and $D$ the \omterm{def:m2:mortgages-and-agencies:convexity}{effective duration}: after a fall, the holder's
DV01 shrinks and it must buy duration (receive fixed or buy bonds); after a
rise, sell it. The trades go in the direction of the move.
\end{proposition}

\begin{proof}
The portfolio's DV01 is $F P D/100 \times 10^{-4}$; a zero-duration holder's
hedges carry the opposite DV01, which must follow any change in it.
\end{proof}

This is the mechanism of the hook. Not every holder hedges, and some buy
options (\cref{ch:m2:the-rates-options-market}) rather than trade dynamically,
but the dealers, servicers and mortgage investors who do hedge dynamically all
trade in the direction the market has just moved, and in size. The Federal
Reserve economists who made the 2003 estimate found that the volatility implied
by swap options rises when the mortgage market's \omterm{def:m2:mortgages-and-agencies:cpr}{prepayment} risk rises, as if
investors expected the hedging to amplify rate moves. They also recalled 1994,
when long rates rose 133 basis points between October 1993 and May 1994 while
the Fed tightened 125, a move attributed at the time to mortgage holders
selling Treasuries as their duration surged.

\begin{omfigure}
\begin{tikzpicture}
\begin{axis}[width=11cm,height=4.8cm,
  ylabel={USD billion},
  tick label style={font=\small},label style={font=\small},
  xmin=2007,xmax=2026.75,ymin=0,ymax=3000,grid=major,grid style={gray!25},
  xtick={2008,2010,2012,2014,2016,2018,2020,2022,2024,2026},xticklabel style={/pgf/number format/1000 sep={}},
  yticklabel style={/pgf/number format/1000 sep={\,}},
  table/col sep=comma]
\addplot[omThm,very thick] table[x=t,y=bn]{figdata/markets-2/12-mortgages-and-agencies/fedmbs.csv};
\end{axis}
\end{tikzpicture}

\omcaption{Mortgage-backed securities held outright by the Federal Reserve,
month ends, 2007 to August 2026. Purchases began in January 2009, resumed in
2012 and in 2020, and the holdings have run off since 2022. Data: FRED series
WSHOMCB, from the Federal Reserve's H.4.1 release.}\label{fig:m2:mortgages-and-agencies:fed}
\end{omfigure}

The central bank's holdings change the picture. A central bank does not
hedge the duration of its holdings as a bank or a servicer does, so its
purchases take convexity hedging out of the market, and its run-off puts it
back. The figure is also a reminder that the largest owner of
the asset can change its mind.

\section{Tutorial: pass-through cash flows and effective duration}\label{tut:m2:mortgages-and-agencies:prepay}

\begin{tutorial}
\textbf{Goal.} Project the monthly cash flows of a \omterm{def:m2:mortgages-and-agencies:mbs}{pass-through} under any
\omterm{def:m2:mortgages-and-agencies:cpr}{prepayment} function, price it, and measure its \omterm{def:m2:mortgages-and-agencies:convexity}{effective duration} and
convexity with a refinancing S-curve; price a \omterm{def:m2:mortgages-and-agencies:tba}{dollar roll}. \textbf{End state:}
\cref{fig:m2:mortgages-and-agencies:psa,fig:m2:mortgages-and-agencies:priceyield},
\cref{ex:m2:mortgages-and-agencies:pool,ex:m2:mortgages-and-agencies:roll} and
the numbers of the weekend problem.
\begin{tutsteps}
\item \textbf{Cash flows}: level payment on the remaining balance and term,
scheduled principal, then \omterm{def:m2:mortgages-and-agencies:cpr}{prepayment} of a fraction SMM of what is left.
\omcode{code/firm/prepay/firm_prepay.py}{38}{51}{Monthly cash flows of a pass-through.}\label{lst:m2:mortgages-and-agencies:flows}
\item \textbf{Speeds and risk}: the S-curve, the \omterm{def:m2:mortgages-and-agencies:convexity}{effective duration} and
convexity by re-running the cash flows at shifted rates, and the roll's
financing rate of \cref{prop:m2:mortgages-and-agencies:roll}.
\omcode{code/firm/prepay/firm_prepay.py}{65}{84}{Refinancing S-curve, effective risk and the dollar roll.}\label{lst:m2:mortgages-and-agencies:risk}
\item \textbf{Run} \texttt{mbs\_demo.risk\_table()},
\texttt{mbs\_demo.convexity\_hedge()}, \texttt{mbs\_demo.roll\_example()}
and \texttt{fig\_mbs.py}.
\end{tutsteps}
\textbf{What to change next.} Flatten the S-curve (a smaller slope) and watch
the \omterm{def:m2:mortgages-and-agencies:convexity}{negative convexity} shrink; season the pool from zero months and see the
ramp delay the refinancing response.
\end{tutorial}

\section{Build: the pass-through engine}\label{bld:m2:mortgages-and-agencies:prepay}

\begin{build}
\bfield{Purpose} The miniature firm makes markets in TBAs, finances them with
rolls and hedges their duration: it needs cash flows under any speed, prices,
weighted average lives, \omterm{def:m2:mortgages-and-agencies:convexity}{effective durations}, and roll financing rates.
\bfield{Interface} \texttt{psa\_cpr(age, speed)}; \texttt{smm(cpr)};
\texttt{level\_payment}; \texttt{cash\_flows(balance, wac, coupon, term, age,
cpr)} returning \texttt{Flow} records; \texttt{price(flows, y, face)};
\texttt{wal}; \texttt{refi\_cpr(incentive, \ldots)};
\texttt{effective\_risk(p, y, dy)}; \texttt{roll\_financing\_rate(front, back,
coupon, paydown, days)}.
\bfield{Rules} Monthly periods; \omterm{def:m2:mortgages-and-agencies:cpr}{prepayment} applied after scheduled principal;
interest to investors at the net coupon on the opening balance; discounting at
a monthly-compounded yield without the payment delay; one static rate path.
\bfield{Acceptance tests} \texttt{code/firm/prepay/tests/}: the PSA ramp's
defining points; with no \omterm{def:m2:mortgages-and-agencies:cpr}{prepayment} a level-payment mortgage fully amortised;
all principal returned and the price at par when the yield equals the coupon,
at any speed; \omterm{def:m2:mortgages-and-agencies:convexity}{negative convexity} with the S-curve and positive with frozen
speeds; a roll with no drop and no paydown finances at the coupon.
\bfield{Stretch} A payment delay; rate paths simulated from a short-rate model
and an option-adjusted spread (chapter 21 and One Quant Book 6); pool-level
speeds by loan size and age; the TBA \omterm{def:m2:bond-futures:ctd}{cheapest-to-deliver} across pools.
\end{build}

\begin{omsources}
\item Ginnie Mae, \emph{Base Offering Circular}, definition of the PSA
prepayment assumption.
\item A.~Fuster, D.~Lucca and J.~Vickery, ``Mortgage-backed securities'',
Federal Reserve Bank of New York Staff Report 1001, 2022.
\item R.~Perli and B.~Sack, ``Does mortgage hedging amplify movements in
long-term interest rates?'', Finance and Economics Discussion Series 2003-49,
Federal Reserve Board.
\item SIFMA, US mortgage-backed securities statistics; FRED series WSHOMCB and
MORTGAGE30US.
\end{omsources}

\section{Exercises}

\begin{exercise}[$\star$]\label{exo:m2:mortgages-and-agencies:1}
Give the CPR at 100\% PSA in a loan's tenth month, at 150\% PSA in its
twentieth, and at 200\% PSA in its fortieth.
\end{exercise}

\begin{exercise}[$\star$]\label{exo:m2:mortgages-and-agencies:2}
Convert a CPR of 6\% into a single monthly mortality. A pool has USD~50
million left after scheduled principal: how much prepays this month?
\end{exercise}

\begin{exercise}[$\star$]\label{exo:m2:mortgages-and-agencies:3}
Name the six parameters of a TBA trade. Which pools will a seller deliver, and
why do some pools trade separately?
\end{exercise}

\begin{exercise}[$\star\star$]\label{exo:m2:mortgages-and-agencies:4}
In the model of \cref{ex:m2:mortgages-and-agencies:roll}, give the roll's
implied financing rate if the drop were zero, and explain why it is close to
the coupon.
\end{exercise}

\begin{exercise}[$\star\star$]\label{exo:m2:mortgages-and-agencies:5}
From \cref{fig:m2:mortgages-and-agencies:priceyield}, why is the
\omterm{def:m2:mortgages-and-agencies:mbs}{pass-through}'s price at 4\% only 102.88 when the frozen-speed pool is worth
110.58?
\end{exercise}

\begin{exercise}[$\star\star$]\label{exo:m2:mortgages-and-agencies:6}
Give the \omterm{def:m2:mortgages-and-agencies:convexity}{effective duration} and convexity of the chapter's \omterm{def:m2:mortgages-and-agencies:mbs}{pass-through} at
5.5\%, 6\% and 6.5\%, and explain where the \omterm{def:m2:mortgages-and-agencies:convexity}{negative convexity} is strongest.
\end{exercise}

\begin{exercise}[$\star\star\star$]\label{exo:m2:mortgages-and-agencies:7}
\emph{Coding.} With \texttt{cash\_flows} give the weighted average life of the
new pool at 0, 50, 100, 200 and 300\% PSA.
\end{exercise}

\begin{exercise}[$\star\star\star$]\label{exo:m2:mortgages-and-agencies:8}
\emph{Find the flaw.} ``Agency mortgages are guaranteed by the government,
so they are as safe as Treasuries and should yield the same.'' Correct it.
\end{exercise}

\section{Problem: The Refinancing Wave}

\begin{problem}[{Weekend problem --- how a rally turns into a bid for duration}]\label{pb:m2:mortgages-and-agencies:1}
An investor holds USD~10 billion face of the chapter's seasoned 6\%
\omterm{def:m2:mortgages-and-agencies:mbs}{pass-through} (6.5\% loans, 36 months old), hedged to zero duration by paying
fixed on ten-year swaps. The discount yield is 6\%, the primary mortgage rate
6.5\%, and speeds follow the chapter's S-curve. Ten-year swaps are priced on a
flat 6\% annual curve.

\medskip\noindent\textbf{Part I --- At the money.}
\begin{enumerate}
\item Give the pool's CPR, price and \omterm{def:m2:mortgages-and-agencies:convexity}{effective duration}.
\item Give the portfolio's DV01.
\item Give the DV01 of USD~100 million of the ten-year swap.
\item What notional of swaps is the investor paying fixed on?
\item Give the effective convexity, and say what its sign means.
\end{enumerate}

\medskip\noindent\textbf{Part II --- The rally.}
\begin{enumerate}[resume]
\item Rates fall 50 basis points. Give the new CPR, price and \omterm{def:m2:mortgages-and-agencies:convexity}{effective duration}.
\item Give the portfolio's new DV01.
\item By how much must the hedge change, and in which direction?
\item What notional of ten-year swaps must be received?
\item Give the \omterm{def:m2:mortgages-and-agencies:convexity}{effective duration} if rates instead rise 50 basis points.
\end{enumerate}

\medskip\noindent\textbf{Part III --- The market.}
\begin{enumerate}[resume]
\item The 2003 estimate for the whole market was USD~289 billion of ten-year
swaps for a 50-basis-point move. What does the investor's trade represent of it?
\item Why does the hedge push rates further in the direction of the move?
\item Who takes the other side?
\item How would buying options instead change the investor's trading?
\item Why did the Federal Reserve's purchases reduce this effect?
\end{enumerate}

\medskip\noindent\textbf{Part IV --- Judgement.}
\begin{enumerate}[resume]
\item Why does the model overstate the precision of the answer?
\item Why does a new pool respond less to the rally?
\item What would a servicer, paid a fee on the outstanding balance, do after the rally?
\item State the \emph{named result}: the convexity hedge for USD~10 billion
after a 50-basis-point rally.
\item In one sentence: why do mortgage holders buy duration after rates fall?
\end{enumerate}
\end{problem}

\section{Interview questions}

\begin{interviewq}[$\star$ \iqroles{trader, researcher}]\label{iq:m2:mortgages-and-agencies:1}
Why does a mortgage-backed security have \omterm{def:m2:mortgages-and-agencies:convexity}{negative convexity}?
\end{interviewq}

\begin{interviewq}[$\star$ \iqroles{trader, bank}]\label{iq:m2:mortgages-and-agencies:2}
What is a TBA, and why does the market trade mostly TBAs rather than pools?
\end{interviewq}

\begin{interviewq}[$\star\star$ \iqroles{trader}]\label{iq:m2:mortgages-and-agencies:3}
Explain a \omterm{def:m2:mortgages-and-agencies:tba}{dollar roll}. When is a roll special, and what do you do if it is?
\end{interviewq}

\begin{interviewq}[$\star\star$ \iqroles{researcher}]\label{iq:m2:mortgages-and-agencies:4}
How would you build a \omterm{def:m2:mortgages-and-agencies:cpr}{prepayment} model, and how would you know it is good?
\end{interviewq}

\begin{interviewq}[$\star\star$ \iqroles{researcher, trader}]\label{iq:m2:mortgages-and-agencies:5}
How do mortgage hedging flows affect the Treasury and swap markets?
\end{interviewq}

\begin{interviewq}[$\star\star\star$ \iqroles{developer, researcher}]\label{iq:m2:mortgages-and-agencies:6}
Design the system that computes the \omterm{def:m2:mortgages-and-agencies:convexity}{effective duration} of a book of a hundred
thousand pools every hour.
\end{interviewq}
